% from ega2-1-fr.tex
\oldpage[II]{96}

et résultent donc aussitôt de (\hyperref[II.5.1.2-fr]{5.1.2}) et
(\hyperref[II.5.1.4-fr]{5.1.4}), compte tenu de ce que
$f_*(\mathcal{F})$ est quasi-cohérent lorsque $f$ est séparé
(\hyperref[I.9.2.2-fr]{I, 9.2.2, a)}).

\begin{corollary}[5.1.7]
\label{II.5.1.7-fr}
Soit $f:X\to Y$ un morphisme quasi-affine. Pour tout ouvert $U$ de $Y$, la
restriction $f^{-1}(U)\to U$ de $f$ est quasi-affine.
\end{corollary}

\begin{corollary}[5.1.8]
\label{II.5.1.8-fr}
Soient $Y$ un schéma affine, $f:X\to Y$ un morphisme quasi-compact. Pour que
$f$ soit quasi-affine, il faut et il suffit que $X$ soit un schéma
quasi-affine.
\end{corollary}

C'est une conséquence immédiate de (\hyperref[II.5.1.6-fr]{5.1.6}) et
(\hyperref[II.4.6.6-fr]{4.6.6}).

\begin{corollary}[5.1.9]
\label{II.5.1.9-fr}
Soient $Y$ un schéma quasi-compact ou un préschéma dont l'espace sous-jacent
est noethérien, $f:X\to Y$ un morphisme de type fini. Si $f$ est
quasi-affine, il existe une sous-$\mathcal{O}_Y$-Algèbre quasi-cohérente
$\mathcal{B}$ de $\mathcal{A}(X)=f_*(\mathcal{O}_X)$, de type fini
(\hyperref[I.9.6.2-fr]{I, 9.6.2}), telle que le morphisme
$X\to\operatorname{Spec}(\mathcal{B})$ correspondant à l'injection canonique
$\mathcal{B}\to\mathcal{A}(X)$ (\hyperref[II.1.2.7-fr]{1.2.7}) soit une
immersion. En outre, toute sous-$\mathcal{O}_Y$-Algèbre quasi-cohérente de
type fini $\mathcal{B}'$ de $\mathcal{A}(X)$, contenant $\mathcal{B}$,
possède la même propriété.
\end{corollary}

En effet, $\mathcal{A}(X)$ est limite inductive de ses
sous-$\mathcal{O}_Y$-Algèbres quasi-cohérentes de type fini
(\hyperref[I.9.6.5-fr]{I, 9.6.5})~; le résultat est alors un cas particulier
de (\hyperref[II.3.8.4-fr]{3.8.4}), compte tenu de l'identification de
$\operatorname{Spec}(\mathcal{A}(X))$ et de
$\operatorname{Proj}(\mathcal{A}(X)[T])$
(\hyperref[II.3.1.7-fr]{3.1.7}).

\begin{proposition}[5.1.10]
\label{II.5.1.10-fr}
\begin{enumerate}
  \item[\rm{(i)}] Un morphisme quasi-compact $X\to Y$ qui est un
  homéomorphisme de l'espace sous-jacent à $X$ sur un sous-espace de
  l'espace sous-jacent à $Y$ (en particulier une immersion fermée) est
  quasi-affine.
  \item[\rm{(ii)}] Le composé de deux morphismes quasi-affines est
  quasi-affine.
  \item[\rm{(iii)}] Si $f:X\to Y$ est un $S$-morphisme quasi-affine,
  $f_{(S')}:X_{(S')}\to Y_{(S')}$ est un morphisme quasi-affine pour toute
  extension $S'\to S$ du préschéma de base.
  \item[\rm{(iv)}] Si $f:X\to Y$ et $g:X'\to Y'$ sont deux $S$-morphismes
  quasi-affines, $f\times_S g$ est quasi-affine.
  \item[\rm{(v)}] Si $f:X\to Y$, $g:Y\to Z$ sont deux morphismes tels que
  $g\circ f$ soit quasi-affine, et si $g$ est séparé ou l'espace sous-jacent
  à $X$ localement noethérien, $f$ est quasi-affine.
  \item[\rm{(vi)}] Si $f$ est un morphisme quasi-affine, il en est de même
  de $f_{\mathrm{red}}$.
\end{enumerate}
\end{proposition}

Compte tenu du critère (\hyperref[II.5.1.6-fr]{5.1.6, c')}), (i), (iii),
(iv), (v) et (vi) découlent aussitôt de
(\hyperref[II.4.6.13-fr]{4.6.13, (i bis), (iii), (iv), (v) et (vi)})
respectivement. Pour démontrer (ii), on peut se borner au cas où $Z$ est
affine, et alors l'assertion résulte directement de
(\hyperref[II.4.6.13-fr]{4.6.13, (ii)}), appliqué à
$\mathcal{L}=\mathcal{O}_X$ et $\mathcal{K}=\mathcal{O}_Y$.

\begin{remark}[5.1.11]
\label{II.5.1.11-fr}
Soient $f:X\to Y$, $g:Y\to Z$ deux morphismes tels que $X\times_ZY$ soit
localement noethérien. Alors l'immersion graphe
$\Gamma_f:X\to X\times_ZY$ est quasi-affine, étant quasi-compacte
(\hyperref[I.6.3.5-fr]{I, 6.3.5}), et
(\hyperref[I.5.5.12-fr]{I, 5.5.12}) montre que, dans (v), la conclusion reste
valable si on supprime l'hypothèse que $g$ est séparé.
\end{remark}

\begin{proposition}[5.1.12]
\label{II.5.1.12-fr}
Soient $f:X\to Y$ un morphisme quasi-compact, $g:X'\to X$ un morphisme
quasi-affine. Si $\mathcal{L}$ est un $\mathcal{O}_X$-Module ample pour
$f$, $g^*(\mathcal{L})$ est un $\mathcal{O}_{X'}$-Module ample pour
$f\circ g$.
\end{proposition}

En effet, comme $\mathcal{O}_{X'}$ est très ample pour $g$, et que la
question est locale sur $Y$ (\hyperref[II.4.6.4-fr]{4.6.4}), il résulte de
(\hyperref[II.4.6.13-fr]{4.6.13, (ii)}) qu'il existe (pour $Y$ affine) un
entier $n$ tel que $g^*(\mathcal{L})^{\otimes n}$ soit ample pour
$f\circ g$, donc $g^*(\mathcal{L})$ est ample pour $f\circ g$
(\hyperref[II.4.6.9-fr]{4.6.9}).

